\documentclass[10pt,a4paper]{amsart}
\usepackage[margin=19mm]{geometry}
\usepackage{amsmath,amssymb,mathtools}
\usepackage[scr=rsfs,cal=euler]{mathalfa}
\usepackage{xfrac}
\usepackage[unicode,hidelinks,bookmarks=false]{hyperref}
\newcommand{\ie}{i.e.\ }
\newcommand{\cf}{cf.\ }
\newcommand{\resp}{respectively\ }
\newcommand{\Subsection}{\subsection}
\providecommand{\nobreakdash}{\nobreak\hbox{-}}
\setlength{\emergencystretch}{2em}
\multlinegap0pt
\usepackage{amssymb}
\usepackage{tikz}
\newtheorem{lemma}{Lemma}
\title{About the Cram\'er Large Deviation Property for Bell Polynomials}
\author{Sophia Li, Shannon Starr}
\date{Source: arXiv:2609.06281v2 (CC BY 4.0)}
\begin{document}
\maketitle

\noindent\emph{Source and licence.} About the Cram\'er Large Deviation Property for Bell Polynomials. Version arXiv:2609.06281v2 (CC BY 4.0), distributed by arXiv under the Creative Commons Attribution 4.0 International licence, http://creativecommons.org/licenses/by/4.0/, as recorded in the arXiv licence metadata for this exact version and shown on the abstract page https://arxiv.org/abs/2609.06281v2. Full licence notice: \texttt{licenses/arxiv-2609.06281-CC-BY-4.0.txt}.\par\smallskip

\noindent\textit{Editorial scope.} Counted unit: the article's lemma lemma (source0.tex lines 242--244). The article's complete original proof of it is delivered with it (source0.tex lines 245--260, one \texttt{proof} environment of 16 lines). No mathematics is written, repaired or re-derived: the statement and the proof are the article's own lines, in the article's own notation, and no retained line is substituted or re-flowed.  No further source-local context is needed: every cross-reference of every retained line resolves inside the two delivered spans.  Nothing was omitted from any retained span, so the package carries no omission record.  No retained line contains a citation, so the wrapper carries no bibliography and no bibliography entry.  No retained line contains a figure environment, an image-inclusion command or drawing code, so no image is delivered and the package declares no asset dependency.  Editorial formatting only, in addition: an AMS \texttt{amsart} preamble that restates the article's own declarations and macros used by the retained lines, namely the environments \texttt{lemma} on separate counters, together with the standard packages those lines need.  The retained environments print in this document as Lemma 1 in order of appearance, whereas the article prints the same items with the numbering of its own full document.  The complete native source of the article as distributed in the arXiv e-print for this exact version is bundled unchanged as \texttt{sources/209515/source0.tex}.
\begin{lemma}
If $R=\infty$ then $\lim_{r \to \infty} \mathcal{M}(r)=\infty$.
\end{lemma}

\begin{proof}
By Chebyshev's inequality, for $0<t<T$ we have
$$
\sum_{n=1}^{N} \mu_n(e^T)\, \leq\, \sum_{n=1}^{N} e^{(N-n)(T-t)} \mu_{n}(e^T)\,
\leq\, e^{N(T-t)}\, \cdot 
\frac{\sum_{n=1}^{N} e^{nt} Q_{1,n}}{\sum_{n=1}^{\infty} e^{nT} Q_{1,n}}\, ,
$$
which is bounded by $\exp\left(N(T-t)+\ln(F(e^t))-\ln(F(e^T))\right)$.
But $F(e^T)>Q_{1,N+1} e^{(N+1)T}$. So $\ln(F(e^T))>(N+1)T+\ln(Q_{1,N+1})$.
Note that we have assumed that $w_{N+1}\geq 1$ so that $\ln(Q_{1,N+1})=\ln(w_{N+1}/(N+1)!)$ is not $-\infty$.
This implies
$$
\lim_{T \to \infty} \sum_{n=1}^{N} \mu_n(e^T)\, =\, 0\, .
$$
So $\lim_{r \to \infty} \mathcal{M}(r)>N+1$. But since $N$ is arbitrary, this proves the lemma.
\end{proof}

\end{document}
